\clearpage
\section{Zentraler Interaktionsablauf}

\subsection{Standardfall}

Der Standardfall beschreibt eine erfolgreiche Bestätigung des vorgeschlagenen
Lieferzeitfensters:

\begin{enumerate}
  \item Der Disponent plant Auftrag und Tour in DISPO.
  \item Er startet die Avisierung für einen einzelnen Auftrag.
  \item DISPO übergibt die geplanten Lieferdaten an den Avisierungsservice.
  \item Der Avisierungsservice versendet eine Nachricht mit persönlichem Link.
  \item Der Kunde öffnet die Bestätigungsseite und bestätigt das Zeitfenster.
  \item Der Avisierungsservice speichert die Rückmeldung als \texttt{CONFIRMED}.
  \item Das Ergebnis erscheint im Avisierungsdashboard und wird an DISPO zurückgemeldet.
\end{enumerate}

\begin{figure}[H]
\centering
\begin{tikzpicture}[font=\tiny,>=Latex]
  \matrix (m) [matrix of nodes,column sep=0.35cm,row sep=0.55cm,nodes={minimum width=1.85cm,minimum height=0.5cm,align=center}] {
    |[draw,fill=SoftGray,rounded corners]| Disponent &
    |[draw,fill=SoftGray,rounded corners]| DISPO &
    |[draw=MinovaBlue,fill=SoftBlue,rounded corners]| Avisierungsservice &
    |[draw,fill=SoftGray,rounded corners]| Nachrichtenkanal &
    |[draw=MinovaBlue,fill=SoftBlue,rounded corners]| Bestätigungsseite &
    |[draw=MinovaBlue,fill=SoftBlue,rounded corners]| Dashboard \\
  };
  \foreach \x in {1,...,6} {
    \draw[dashed,gray] (m-1-\x.south) -- ++(0,-7.5cm);
  }

  \coordinate (y1) at ([yshift=-0.9cm]m-1-1.south);
  \coordinate (y2) at ([yshift=-1.8cm]m-1-1.south);
  \coordinate (y3) at ([yshift=-2.7cm]m-1-1.south);
  \coordinate (y4) at ([yshift=-3.6cm]m-1-1.south);
  \coordinate (y5) at ([yshift=-4.5cm]m-1-1.south);
  \coordinate (y6) at ([yshift=-5.4cm]m-1-1.south);
  \coordinate (y7) at ([yshift=-6.3cm]m-1-1.south);
  \coordinate (y8) at ([yshift=-7.2cm]m-1-1.south);

  \draw[->,thick] (m-1-1.center |- y1) -- node[above]{Avisierung starten} (m-1-2.center |- y1);
  \draw[->,thick] (m-1-2.center |- y2) -- node[above]{Lieferdaten} (m-1-3.center |- y2);
  \draw[->,thick] (m-1-3.center |- y3) -- node[above]{Link versenden} (m-1-4.center |- y3);
  \draw[->,thick] (m-1-4.center |- y4) -- node[above]{Nachricht} (m-1-5.center |- y4);
  \draw[->,thick] (m-1-5.center |- y5) -- node[above]{Lieferdaten laden} (m-1-3.center |- y5);
  \draw[->,thick] (m-1-5.center |- y6) -- node[above]{bestätigen} (m-1-3.center |- y6);
  \draw[->,thick] (m-1-3.center |- y7) -- node[above]{Status anzeigen} (m-1-6.center |- y7);
  \draw[->,thick] (m-1-3.center |- y8) -- node[above]{Ergebnis melden} (m-1-2.center |- y8);
\end{tikzpicture}
\caption{Interaktion im Standardfall}
\label{fig:sequence}
\end{figure}


\subsection{Fälle mit Handlungsbedarf}

Lehnt der Kunde das Zeitfenster ab, wird der Fall als \texttt{REJECTED} mit seinem Kommentar
angezeigt. Bleibt eine Antwort innerhalb der vorgesehenen Frist aus, wird er als
\texttt{NO\_RESPONSE} sichtbar. Beide Fälle erfordern eine Entscheidung des Disponenten.

Der Disponent kann die unveränderte Anfrage über einen anderen Kanal erneut versenden. Muss
das Lieferzeitfenster geändert werden, erfolgt diese Änderung zuerst in DISPO; anschlie\ss{}end
wird eine neue Avisierung ausgelöst. Wird die Situation telefonisch geklärt, kann der
Disponent den Fall in der Detailansicht manuell bestätigen. Diese Aktion wird mit Benutzer
und Zeitpunkt protokolliert und ebenfalls an DISPO zurückgemeldet.
